\subsection*{A fixed finite-floor experiment and its cutoff cost}

Proposition~\ref{prop:v136-bounded-floor} permits a fixed ordinary
negative error of any finite size. This does make some finite-window
certificates cheaper. It does not by itself remove the scalar-tail
cost in Proposition~\ref{prop:v125-cutoff-cost}.

\begin{proposition}[Additive shifts and the remaining cutoff barrier]
\label{prop:v138-shifted-floor}
In the orthonormal Fourier splitting of
Proposition~\ref{prop:v116-directional-schur}, let $\delta\ge0$ and
suppose $T\succeq D_g$, with increasing $g_n$ and
$g_{N+1}+\delta>0$. For a fixed finite solve $Z$ set
\[
 K_\delta=K_Z+\delta(I+Z^*Z),\qquad
 R_\delta=R-\delta Z.
\]
If $U_J$ encloses the remote residual beyond the support of $Z$, then
\begin{equation}\label{eq:v138-shifted-lower}
 \mathcal L_\delta:=K_\delta-
 \sum_{N<n\le J}\frac{R_{\delta,n}^*R_{\delta,n}}{g_n+\delta}
 -\frac{U_J}{g_{J+1}+\delta}\succeq0
 \quad\Longrightarrow\quad
 \mathsf W_\lambda\succeq-\delta I.
\end{equation}
The displayed sum is in one normalized parity basis. Verifying both
parities proves the inequality on the full complex form domain.
For the present scalar comparison with $c=0$, its prerequisite forces
\begin{equation}\label{eq:v138-shifted-cost}
 N+1>L\exp(M_{\phi,\lambda}-\delta),\qquad L=2\log\lambda.
\end{equation}
Consequently any bounded choice of $\delta$ retains
$\liminf_{\lambda\to\infty}\log((N+1)/L)/\lambda\ge1$.
A polynomial cutoff in this particular comparison would instead
require $\delta\ge M_{\phi,\lambda}-O(\log\lambda)$.
\end{proposition}
\begin{proof}
Apply the verified-solve identity to $\mathsf W_\lambda+\delta I$.
Its off-diagonal block is still $B$, its tail is $T+\delta I$, and
expansion gives the stated $K_\delta$ and $R_\delta$. Inverse order
gives $(T+\delta I)^{-1}\preceq(D_g+\delta I)^{-1}$.
The shifted remote rows beyond the solve support are unchanged.
The moment enclosure and completion of the square prove
\eqref{eq:v138-shifted-lower} on the closed form domain.
For the cost, positivity of the first scalar weight implies
$\log((N+1)/L)>M_{\phi,\lambda}-\delta$, since its other losses
are nonnegative. Use
$M_{\phi,\lambda}\ge\|\mathsf T_{{\rm pr},\lambda}\|$ and
Proposition~\ref{prop:v119-prime-scale}. This is a limitation of
the chosen scalar comparison, not of every matrix tail metric.
\end{proof}

Fresh outward evaluation of all window-dependent constants gives the
following smallest admitted integer $N\ge1$ with $g_{N+1}+\delta>0$:
\begin{center}
\begin{tabular}{rrrrr}
\toprule
$\lambda$ & $\delta=10^{-3}$ even & $\delta=10^{-3}$ odd
 & $\delta=8$ even & $\delta=8$ odd\\
\midrule
5&1501&7217&1&4\\
6&7481&35985&3&14\\
8&124317&598024&43&204\\
\bottomrule
\end{tabular}
\end{center}
Each threshold except the lower endpoint $N=1$ has a certified
negative predecessor. These are only remote-tail gates. They show
why a moderately sized constant, rather than $10^{-3}$, is the useful
finite experiment suggested by the bounded-floor reduction.

\begin{proposition}[Complete finite floors at three windows]
\label{prop:v138-three-floors}
For each $\lambda\in\{5,6,8\}$ and both reflection sectors,
\begin{equation}\label{eq:v138-three-floors}
 \mathsf W_\lambda\succeq-8I,\qquad
 -8\le\inf\sigma(\mathsf W_\lambda^{\pm})<10^{-16}.
\end{equation}
The lower estimate includes the complete infinite Fourier tail;
the upper estimate comes from a certified finite-support trial.
By physical support consistency, the lower estimate also holds for
every $1<\lambda\le8$. No cofinal lower estimate is asserted.
\end{proposition}
\begin{proof}[Complete interval certificate]
Use the canonical form and logarithmic domain of
Proposition~\ref{prop:v114-weil-core}, with ordinary physical
$L^2$ normalization and its established Fourier form core. Freeze
the same parameters at all three windows:
\[
 \delta=8,\quad N=256,\quad J=4096,\quad r=16,
 \quad\tau=1/10,\quad Z=0.
\]
The even head is $0{:}256$, the odd head $1{:}256$. There is no
conjugate-gradient solve or positive $10^{-8}D$ tail prerequisite.
In the remote-moment formula take $M=N$, so $G$ is the head embedding
and $G^*G=I$. The moment proof only requires $M<J$; a nonempty
intermediate solve block is unnecessary when $Z=0$.
Here $K_8=F+8I$, $R_8=B$, and all actual rows through $J$ are retained.
The entire remainder beyond $J$ is bounded by
Proposition~\ref{prop:v116-remote-moment}, including its leading
moment Gram and geometric remainder.

At each window recompute $L$, every strict-cutoff prime power,
the separate logarithmic diagonal, $M_{\phi,\lambda}$, the parity
tail constants and $B_*$. The coefficient enclosure retains $64$
archimedean series terms and the analytic infinite-series remainder.
The first runs at $160$ bits freeze inverse-Cholesky congruences as
exact IEEE dyadics. A fresh $256$-bit evaluation, using the identical
dyadics, verifies every row of
$C\mathcal L_8C^*$ is greater than $999/1000$ after subtraction
of the absolute off-diagonal row sum. The matrices $C$ are triangular
with nonzero diagonal. Thus every complete lower matrix is positive.
The code also checks exact finite-support Rayleigh quotients: each
is strictly positive and less than $10^{-16}$. These are upper
bounds for the true lower spectral edge, not signs of that edge.
The two parities give the full complex-space inequality.
Finally, zero extension of a test from a smaller physical interval
preserves its ordinary norm and Weil form. Apply the $\lambda=8$
bound and use the form core to obtain the assertion for $\lambda\le8$.
\end{proof}

The near-one congruence row margins are not invariant headroom.
For a comparable certificate quantity put
$\mathcal U_8=K_8-\mathcal L_8$ and
\[
 \mathfrak h_\lambda^\pm
 =1-\lambda_{\max}(\mathcal U_8,K_8).
\]
The same exact congruence verifies
$C\mathcal L_8C^*\succeq(999/1000)I$ and
$CK_8C^*\preceq M I$. Hence
$\mathfrak h_\lambda^\pm\ge(999/1000)/M$.
The following are downward-rounded lower bounds for this generalized
margin of the \emph{shifted} certificate:
\begin{center}
\begin{tabular}{rrrrr}
\toprule
$\lambda$ & even $g_{257}+8$ & odd $g_{257}+8$
 & even $\mathfrak h$ lower & odd $\mathfrak h$ lower\\
\midrule
5&6.2297&4.6549&0.8430&0.9190\\
6&4.6225&3.0457&0.7166&0.7620\\
8&1.8111&0.2302&0.2474&0.1022\\
\bottomrule
\end{tabular}
\end{center}
These margins quantify enclosure headroom with the same target and
construction. Their decrease is not evidence of a negative physical
eigenvalue; the actual lower edge remains bracketed only between
$-8$ and the small positive trial quotient. A worsening lower bound
cannot establish the divergence in
\eqref{eq:v136-negative-divergence}. A negative Rayleigh upper bound
would have a different logical role, and none was found here.

This is a continuation of the existing Schur/moment and shifted-inertia
routes, not a new arithmetic positivity mechanism. General certified
finite-window lower and trial upper bounds also appear in
\cite{Zhu2026}; no external numerical certificate is imported.
The actual unproved target remains one finite ordinary lower constant
on a cofinal family, on the full physical form or on a complement
meeting the residual hypotheses of Proposition~\ref{prop:v136-bounded-floor}.
The present bound covers every physical subspace at the listed windows
without identifying $C_a^{\mathrm{low}}$ with a rank-one complement.
All endpoint, scale and sampler graph-defect obligations are unchanged.
No growing-window G2 gap is closed.

